% Originally: content/week-12.tex, \section{Solutions}

\section{Solutions}
\label{sec:pullbacks_solutions}

\begin{exercisesolution}[Solution of \cref{ex:12.9}]
\textbf{(a)} Writing each $d\psi^j$ in the coordinate basis, $d\psi^j = \sum_{i=1}^n
\frac{\partial\psi^j}{\partial x^i}\,dx^i$, so
\[\omega = \left(\sum_{i_1=1}^n \frac{\partial\psi^1}{\partial x^{i_1}}dx^{i_1}\right)\wedge
  \cdots\wedge\left(\sum_{i_n=1}^n \frac{\partial\psi^n}{\partial x^{i_n}}dx^{i_n}\right)
  = \sum_{i_1,\ldots,i_n} \frac{\partial\psi^1}{\partial x^{i_1}}\cdots
  \frac{\partial\psi^n}{\partial x^{i_n}}\,dx^{i_1}\wedge\cdots\wedge dx^{i_n}\]
by bilinearity of $\wedge$ (property (1) of \cref{def:wedge_product}). Any term with a
repeated index vanishes, since $dx^i\wedge dx^i=0$ by antisymmetry, so only tuples
$(i_1,\ldots,i_n) = (\sigma(1),\ldots,\sigma(n))$ for a permutation $\sigma \in S_n$
contribute; using $dx^{\sigma(1)}\wedge\cdots\wedge dx^{\sigma(n)} = \sgn(\sigma)\,dx^1\wedge
\cdots\wedge dx^n$,
\[\omega = \left(\sum_{\sigma\in S_n} \sgn(\sigma)\,\frac{\partial\psi^1}{\partial
  x^{\sigma(1)}}\cdots\frac{\partial\psi^n}{\partial x^{\sigma(n)}}\right) dx^1\wedge\cdots
  \wedge dx^n.\]
The coefficient is exactly the Leibniz expansion of $\det(D\psi)$, so
\[\boxed{\ d\psi^1\wedge\cdots\wedge d\psi^n = \det(D\psi)\,dx^1\wedge\cdots\wedge dx^n.\ }\]

\textbf{(b)} Let $\eta = f(y)\,dy^1\wedge\cdots\wedge dy^n$ be an $n$-form on $V$. Its
pullback under $\psi$ is $\psi^*\eta = f(\psi(x))\,d\psi^1\wedge\cdots\wedge d\psi^n$, which by
part (a) equals $f(\psi(x))\det(D\psi(x))\,dx^1\wedge\cdots\wedge dx^n$. By
\cref{def:integral_of_form_over_submanifold} (applied with $U$ itself as the $n$-dimensional
submanifold),
\[\int_U \psi^*\eta = \int_U f(\psi(x))\det(D\psi(x))\,dx.\]
Since $\det(D\psi) = \lvert\det(D\psi)\rvert>0$ by hypothesis, the right-hand side is exactly
the ordinary change-of-variables formula, so
\[\int_U \psi^*\eta = \int_U f(\psi(x))\,\lvert\det(D\psi(x))\rvert\,dx = \int_V f(y)\,dy
  = \int_V \eta.\]
Thus $\boxed{\int_U \psi^*\eta = \int_V \eta}$: the change-of-variables formula is exactly the
statement that integration of forms is invariant under pullback by an orientation-preserving
diffeomorphism.
\end{exercisesolution}

% Verified: solved independently, both parts match Sol12, including pullback-invariance.

\begin{exercisesolution}[Solution to \cref{ex:ai_pullback_area_form}]
% Generator: Claude Opus 5 (High)
\textbf{(a)} The pullback of a coordinate differential is the differential of the corresponding
component, so from $x = r\cos\theta$ and $y = r\sin\theta$,
\[\Phi^*(dx) = \cos\theta\,dr - r\sin\theta\,d\theta,
  \qquad
  \Phi^*(dy) = \sin\theta\,dr + r\cos\theta\,d\theta.\]
Wedging, and using $dr\wedge dr = d\theta\wedge d\theta = 0$ together with
$d\theta\wedge dr = -\,dr\wedge d\theta$,
\[\Phi^*(dx\wedge dy)
  = \big(\cos\theta\cdot r\cos\theta\big)\,dr\wedge d\theta
    + \big(-r\sin\theta\cdot\sin\theta\big)\,d\theta\wedge dr
  = r\big(\cos^2\theta + \sin^2\theta\big)\,dr\wedge d\theta
  = r\,dr\wedge d\theta.\]
The factor is exactly the $r$ tabulated for polar coordinates in
\cref{rem:coordinate_jacobians_table}, obtained here without computing a determinant: the
antisymmetry of the wedge did the cofactor expansion.

\textbf{(b)} They measure different things, and the absolute value is the difference.
\Cref{thm:change_of_variables} is a statement about \emph{volume}, which is unsigned: a region
and its mirror image have the same area, so only $\lvert\det\Jac\Phi\rvert$ can appear. The
pullback of a top-degree form is a statement about \emph{oriented} volume, which carries a sign,
and it records $\det\Jac\Phi$ itself.

The swap $\Phi(u_1,u_2) = (u_2,u_1)$ separates them. Its Jacobian is
$\left(\begin{smallmatrix}0&1\\1&0\end{smallmatrix}\right)$ with determinant $-1$, so
\[\Phi^*(dx\wedge dy) = du_2\wedge du_1 = -\,du_1\wedge du_2,\]
while $\lvert\det\Jac\Phi\rvert = 1$ and the change of variables formula does not change at all.
The form notices the reversal; the volume cannot.

\textbf{(c)} \Cref{ex:12.9} concludes that $\int_U\psi^*\eta = \int_V\eta$, with no sign. Drop
the hypothesis $\det > 0$ and take $\psi$ to be the swap above: the pullback acquires a factor
$-1$, so the two integrals differ by a sign and the conclusion is false as stated. The honest
statement without the hypothesis carries $\sgn(\det\Jac\psi)$.

\Cref{thm:change_of_variables} is untroubled by the same map because it never claimed anything
about orientation. This is the precise sense in which the positivity hypothesis in
\cref{ex:12.9} is not a technical convenience: it is the price of upgrading an unsigned identity
to a signed one.
\end{exercisesolution}

\begin{exercisesolution}[Solution to \cref{ex:ai_area_form_over_hemisphere}]
% Generator: Claude Opus 5 (High)
\textbf{(a)} Along $\phi$ the first two coordinates are the parameters themselves, $x = u$ and
$y = v$, so $\phi^*(dx) = du$ and $\phi^*(dy) = dv$, giving
\[\phi^*(dx\wedge dy) = du\wedge dv.\]
The third coordinate never enters, which is the whole reason this integral is easy. It remains to
check that $\phi$ induces the outward orientation: for a graph $z = g(u,v)$,
\[\partial_u\phi\times\partial_v\phi = (-g_u,\ -g_v,\ 1),\]
whose third component is $+1$, so the induced normal points \emph{upward}, and on the upper
hemisphere upward is outward. Hence
\[\int_{S^+} dx\wedge dy = \int_D du\wedge dv = \vol_2(D) = \pi.\]

\textbf{(b)} The pullback computation is identical, since $\psi$ also has $x = u$ and $y = v$,
so again $\psi^*(dx\wedge dy) = du\wedge dv$. But $\psi$ is still a graph, so its induced normal
still points upward, and on the \emph{lower} hemisphere upward is \emph{inward}. So $\psi$ is
orientation-reversing for the prescribed outward orientation, and
\cref{def:integral_of_form_over_submanifold} requires a parametrization compatible with the
chosen orientation. Correcting by the sign,
\[\int_{S^-} dx\wedge dy = -\int_D du\wedge dv = -\pi.\]

\textbf{(c)} The equator is a null set, so the two halves add:
\[\int_{S^2} dx\wedge dy = \pi + (-\pi) = 0.\]
For the second route, $d(x\,dy) = dx\wedge dy$, and $S^2$ is a compact orientable $2$-manifold
with $\partial S^2 = \emptyset$, so \cref{thm:stokes_general} gives
\[\int_{S^2} dx\wedge dy = \int_{S^2} d(x\,dy) = \int_{\partial S^2} x\,dy = 0.\]

\begin{remark}[$dx\wedge dy$ measures the shadow]
\label{rem:area_form_measures_the_shadow}
Part \textbf{(a)} says something worth keeping: integrating $dx\wedge dy$ over an oriented
surface in $\mathbb{R}^3$ returns the \emph{signed} area of its shadow on the $xy$-plane. The
hemisphere has area $2\pi$, but its shadow is the unit disc, and $\pi$ is what the computation
returns. A $2$-form does not measure surface area; the Gram determinant of
\cref{ch:gram_determinant} does that, and the two answers differ by exactly the tilt of the
surface.

Read that way, part \textbf{(c)} is obvious before any computation. The sphere's two halves cast
the same shadow, and a closed surface presents each patch of shadow once from above and once from
below, so everything cancels. That is the geometric content of $\int_M d\omega = 0$ on a closed
$M$, and \cref{q:quiz_q14} draws the other standard consequence from it.
\end{remark}
\end{exercisesolution}
